\documentclass[conference]{IEEEtran}
\IEEEoverridecommandlockouts
% The preceding line is only needed to identify funding in the first footnote. If that is unneeded, please comment it out.
%Template version as of 6/27/2024

\usepackage{cite}
\usepackage{amsmath,amssymb,amsfonts}
\usepackage{algorithmic}
\usepackage{graphicx}
\usepackage{textcomp}
\usepackage{xcolor}
\def\BibTeX{{\rm B\kern-.05em{\sc i\kern-.025em b}\kern-.08em
    T\kern-.1667em\lower.7ex\hbox{E}\kern-.125emX}}
\begin{document}

\title{Benchmarking Quantum vs Classical Approaches for Time-Dependent TSP}

\author{\IEEEauthorblockN{1\textsuperscript{st} Muhammad Ali Khan}
\IEEEauthorblockA{\textit{SuperQ Quantum Computing Inc.} \\
Calgary, Canada \\
makhan@superq.co}
\and
\IEEEauthorblockN{2\textsuperscript{nd} Krishna Ganesh}
\IEEEauthorblockA{\textit{SuperQ Quantum Computing Inc.} \\
Dubai, United Arab Emirates \\
kganesh@superq.co}
\and
\IEEEauthorblockN{3\textsuperscript{rd} Shahadat Hossain}
\IEEEauthorblockA{\textit{Department of Computer Science} \\
\textit{University of Northern British Columbia} \\
Prince George, Canada}
}

\maketitle

\begin{abstract}
The Time-Dependent Traveling Salesperson Problem (TD-TSP) extends the classical TSP by incorporating temporal variability in edge costs, making it highly relevant for real-world applications such as logistics, traffic-aware routing, and dynamic supply chain optimization. However, the time-dependence increases the computational complexity of an already NP-hard standard TSP. This paper investigates potential quantum and quantum-inspired advantage for solving TD-TSP.

We model and solve TD-TSP using (i) quantum annealing on D-Wave hybrid CQM solver, (ii) QAOA implementation on QuEra machine, (iii) quantum-inspired Ising machine by Quanfluence, and (iv) classical Gurobi MILP solver. The modelling and algorithm development takes place on Super platform - autonomous optimization engine developed by SuperQ Quantum that orchestrates optimization workflows across a range of classical and quantum hardware. Super models TD-TSP and generates the initial code base. The authors refined the code through a series of experiments. Both the autonomously generated and human refined code will be made available.

The paper will present optimal values, solver times and communication overhead times from each approach. It will review the state of prior art to benchmark performance. Furthermore, readers will be able to submit, through a web app, instances of TD-TSP to be solved by the authors' code for reproducibility.
\end{abstract}

\begin{IEEEkeywords}
Time-dependent Travelling Salesperson Problem, TSP, Quantum Annealing, QAOA, Ising Machine
\end{IEEEkeywords}

\section{Introduction}
% TODO: Write introduction content
% - Introduce the classical TSP and its importance
% - Explain the Time-Dependent TSP (TD-TSP) extension
% - Discuss the computational challenges (NP-hard complexity)
% - Motivate the exploration of quantum approaches
% - Outline the paper structure

[Introduction placeholder text to be written.]

\section{Literature Review}
% TODO: Write literature review content
% - Prior work on classical TSP algorithms
% - Time-Dependent TSP formulations and solutions
% - Quantum approaches to combinatorial optimization
% - D-Wave quantum annealing for TSP
% - QAOA implementations for routing problems
% - Quantum-inspired algorithms and Ising machines

[Literature review placeholder text to be written.]

\section{Methodology}
Before you begin to format your paper, first write and save the content as a 
separate text file. Complete all content and organizational editing before 
formatting. Please note sections \ref{AA} to \ref{FAT} below for more information on 
proofreading, spelling and grammar.

Keep your text and graphic files separate until after the text has been 
formatted and styled. Do not number text heads---{\LaTeX} will do that 
for you.

\subsection{Abbreviations and Acronyms}\label{AA}
Define abbreviations and acronyms the first time they are used in the text, 
even after they have been defined in the abstract. Abbreviations such as 
IEEE, SI, MKS, CGS, ac, dc, and rms do not have to be defined. Do not use 
abbreviations in the title or heads unless they are unavoidable.

\subsection{Units}
\begin{itemize}
\item Use either SI (MKS) or CGS as primary units. (SI units are encouraged.) English units may be used as secondary units (in parentheses). An exception would be the use of English units as identifiers in trade, such as ``3.5-inch disk drive''.
\item Avoid combining SI and CGS units, such as current in amperes and magnetic field in oersteds. This often leads to confusion because equations do not balance dimensionally. If you must use mixed units, clearly state the units for each quantity that you use in an equation.
\item Do not mix complete spellings and abbreviations of units: ``Wb/m\textsuperscript{2}'' or ``webers per square meter'', not ``webers/m\textsuperscript{2}''. Spell out units when they appear in text: ``. . . a few henries'', not ``. . . a few H''.
\item Use a zero before decimal points: ``0.25'', not ``.25''. Use ``cm\textsuperscript{3}'', not ``cc''.)
\end{itemize}

\subsection{Equations}
Number equations consecutively. To make your 
equations more compact, you may use the solidus (~/~), the exp function, or 
appropriate exponents. Italicize Roman symbols for quantities and variables, 
but not Greek symbols. Use a long dash rather than a hyphen for a minus 
sign. Punctuate equations with commas or periods when they are part of a 
sentence, as in:
\begin{equation}
a+b=\gamma\label{eq}
\end{equation}

Be sure that the 
symbols in your equation have been defined before or immediately following 
the equation. Use ``\eqref{eq}'', not ``Eq.~\eqref{eq}'' or ``equation \eqref{eq}'', except at 
the beginning of a sentence: ``Equation \eqref{eq} is . . .''

\subsection{\LaTeX-Specific Advice}

Please use ``soft'' (e.g., \verb|\eqref{Eq}|) cross references instead
of ``hard'' references (e.g., \verb|(1)|). That will make it possible
to combine sections, add equations, or change the order of figures or
citations without having to go through the file line by line.

Please don't use the \verb|{eqnarray}| equation environment. Use
\verb|{align}| or \verb|{IEEEeqnarray}| instead. The \verb|{eqnarray}|
environment leaves unsightly spaces around relation symbols.

Please note that the \verb|{subequations}| environment in {\LaTeX}
will increment the main equation counter even when there are no
equation numbers displayed. If you forget that, you might write an
article in which the equation numbers skip from (17) to (20), causing
the copy editors to wonder if you've discovered a new method of
counting.

{\BibTeX} does not work by magic. It doesn't get the bibliographic
data from thin air but from .bib files. If you use {\BibTeX} to produce a
bibliography you must send the .bib files. 

{\LaTeX} can't read your mind. If you assign the same label to a
subsubsection and a table, you might find that Table I has been cross
referenced as Table IV-B3. 

{\LaTeX} does not have precognitive abilities. If you put a
\verb|\label| command before the command that updates the counter it's
supposed to be using, the label will pick up the last counter to be
cross referenced instead. In particular, a \verb|\label| command
should not go before the caption of a figure or a table.

Do not use \verb|\nonumber| inside the \verb|{array}| environment. It
will not stop equation numbers inside \verb|{array}| (there won't be
any anyway) and it might stop a wanted equation number in the
surrounding equation.

\subsection{Some Common Mistakes}\label{SCM}
\begin{itemize}
\item The word ``data'' is plural, not singular.
\item The subscript for the permeability of vacuum $\mu_{0}$, and other common scientific constants, is zero with subscript formatting, not a lowercase letter ``o''.
\item In American English, commas, semicolons, periods, question and exclamation marks are located within quotation marks only when a complete thought or name is cited, such as a title or full quotation. When quotation marks are used, instead of a bold or italic typeface, to highlight a word or phrase, punctuation should appear outside of the quotation marks. A parenthetical phrase or statement at the end of a sentence is punctuated outside of the closing parenthesis (like this). (A parenthetical sentence is punctuated within the parentheses.)
\item A graph within a graph is an ``inset'', not an ``insert''. The word alternatively is preferred to the word ``alternately'' (unless you really mean something that alternates).
\item Do not use the word ``essentially'' to mean ``approximately'' or ``effectively''.
\item In your paper title, if the words ``that uses'' can accurately replace the word ``using'', capitalize the ``u''; if not, keep using lower-cased.
\item Be aware of the different meanings of the homophones ``affect'' and ``effect'', ``complement'' and ``compliment'', ``discreet'' and ``discrete'', ``principal'' and ``principle''.
\item Do not confuse ``imply'' and ``infer''.
\item The prefix ``non'' is not a word; it should be joined to the word it modifies, usually without a hyphen.
\item There is no period after the ``et'' in the Latin abbreviation ``et al.''.
\item The abbreviation ``i.e.'' means ``that is'', and the abbreviation ``e.g.'' means ``for example''.
\end{itemize}
An excellent style manual for science writers is \cite{b7}.

\subsection{Authors and Affiliations}\label{AAA}
\textbf{The class file is designed for, but not limited to, six authors.} A 
minimum of one author is required for all conference articles. Author names 
should be listed starting from left to right and then moving down to the 
next line. This is the author sequence that will be used in future citations 
and by indexing services. Names should not be listed in columns nor group by 
affiliation. Please keep your affiliations as succinct as possible (for 
example, do not differentiate among departments of the same organization).

\subsection{Identify the Headings}\label{ITH}
Headings, or heads, are organizational devices that guide the reader through 
your paper. There are two types: component heads and text heads.

Component heads identify the different components of your paper and are not 
topically subordinate to each other. Examples include Acknowledgments and 
References and, for these, the correct style to use is ``Heading 5''. Use 
``figure caption'' for your Figure captions, and ``table head'' for your 
table title. Run-in heads, such as ``Abstract'', will require you to apply a 
style (in this case, italic) in addition to the style provided by the drop 
down menu to differentiate the head from the text.

Text heads organize the topics on a relational, hierarchical basis. For 
example, the paper title is the primary text head because all subsequent 
material relates and elaborates on this one topic. If there are two or more 
sub-topics, the next level head (uppercase Roman numerals) should be used 
and, conversely, if there are not at least two sub-topics, then no subheads 
should be introduced.

\subsection{Figures and Tables}\label{FAT}
\paragraph{Positioning Figures and Tables} Place figures and tables at the top and 
bottom of columns. Avoid placing them in the middle of columns. Large 
figures and tables may span across both columns. Figure captions should be 
below the figures; table heads should appear above the tables. Insert 
figures and tables after they are cited in the text. Use the abbreviation 
``Fig.~\ref{fig}'', even at the beginning of a sentence.

\begin{table}[htbp]
\caption{Table Type Styles}
\begin{center}
\begin{tabular}{|c|c|c|c|}
\hline
\textbf{Table}&\multicolumn{3}{|c|}{\textbf{Table Column Head}} \\
\cline{2-4} 
\textbf{Head} & \textbf{\textit{Table column subhead}}& \textbf{\textit{Subhead}}& \textbf{\textit{Subhead}} \\
\hline
copy& More table copy$^{\mathrm{a}}$& &  \\
\hline
\multicolumn{4}{l}{$^{\mathrm{a}}$Sample of a Table footnote.}
\end{tabular}
\label{tab1}
\end{center}
\end{table}

\begin{figure}[htbp]
\centerline{\includegraphics{fig1.png}}
\caption{Example of a figure caption.}
\label{fig}
\end{figure}

Figure Labels: Use 8 point Times New Roman for Figure labels. Use words 
rather than symbols or abbreviations when writing Figure axis labels to 
avoid confusing the reader. As an example, write the quantity 
``Magnetization'', or ``Magnetization, M'', not just ``M''. If including 
units in the label, present them within parentheses. Do not label axes only 
with units. In the example, write ``Magnetization (A/m)'' or ``Magnetization 
\{A[m(1)]\}'', not just ``A/m''. Do not label axes with a ratio of 
quantities and units. For example, write ``Temperature (K)'', not 
``Temperature/K''.

\section{Experimental Framework}

This section describes the experimental setup, problem formulation, and computational platforms employed in our benchmarking study.

\subsection{Algorithm Development Process}

The initial algorithmic framework for this study was generated using Super, an autonomous optimization engine developed by SuperQ Quantum Computing. Super is designed to orchestrate optimization workflows across heterogeneous computing infrastructure, including classical solvers, quantum annealers, gate-based quantum processors, and quantum-inspired hardware. Given a high-level problem specification, Super automatically generates baseline solver implementations tailored to each target platform.

For this study, Super produced initial implementations for the Time-Dependent Traveling Salesperson Problem (TD-TSP) targeting four distinct computational backends: Gurobi MILP solver, D-Wave quantum annealer, QAOA on gate-based quantum hardware, and the Quanfluence Ising machine. The authors subsequently refined these implementations through extensive experimentation, incorporating problem-specific optimizations such as cluster-based decomposition strategies and local search post-processing. Both the autonomously generated baseline code and the human-refined implementations are made publicly available for reproducibility.

\subsection{Problem Instance Generation}

To ensure fair comparison across all solvers, we employ a standardized methodology for generating TSP instances. For each problem size $n \in \{5, 10, 25, 50, 100\}$, city coordinates are generated uniformly at random within a $100 \times 100$ unit square using a fixed random seed (seed = 42) for reproducibility.

The symmetric distance matrix $D$ is computed using Euclidean distances:
\begin{equation}
D_{ij} = \sqrt{(x_i - x_j)^2 + (y_i - y_j)^2}
\label{eq:euclidean}
\end{equation}
where $(x_i, y_i)$ and $(x_j, y_j)$ represent the coordinates of cities $i$ and $j$, respectively. The diagonal elements $D_{ii} = 0$ for all $i$, and the matrix satisfies the triangle inequality, placing our instances in the metric TSP category.

For Time-Dependent TSP experiments, we extend this formulation by introducing time-varying edge costs. The distance between cities $i$ and $j$ at time slot $t$ is modeled as:
\begin{equation}
D_{ij}^{(t)} = D_{ij} \cdot \tau(t)
\label{eq:td_distance}
\end{equation}
where $\tau(t)$ is a traffic multiplier that varies across predefined time windows (e.g., peak hours: $\tau = 1.5$, off-peak: $\tau = 1.0$, night: $\tau = 0.8$).

\subsection{QUBO Formulation for Quantum Solvers}

For quantum and quantum-inspired solvers, the TSP is encoded as a Quadratic Unconstrained Binary Optimization (QUBO) problem. We introduce $n^2$ binary decision variables $x_{i,t} \in \{0,1\}$, where $x_{i,t} = 1$ indicates that city $i$ is visited at position $t$ in the tour.

The objective function minimizes the total tour distance:
\begin{equation}
C_{\text{obj}} = \sum_{t=0}^{n-1} \sum_{i=0}^{n-1} \sum_{j=0}^{n-1} D_{ij} \cdot x_{i,t} \cdot x_{j,(t+1) \mod n}
\label{eq:qubo_obj}
\end{equation}

Two sets of constraints ensure solution validity. Each city must be visited exactly once:
\begin{equation}
C_{\text{city}} = P \sum_{i=0}^{n-1} \left(1 - \sum_{t=0}^{n-1} x_{i,t}\right)^2
\label{eq:city_constraint}
\end{equation}

Each position must contain exactly one city:
\begin{equation}
C_{\text{pos}} = P \sum_{t=0}^{n-1} \left(1 - \sum_{i=0}^{n-1} x_{i,t}\right)^2
\label{eq:pos_constraint}
\end{equation}

The complete QUBO cost function is $C = C_{\text{obj}} + C_{\text{city}} + C_{\text{pos}}$, where $P$ is a penalty coefficient chosen sufficiently large to ensure constraint satisfaction (typically $P > \max(D_{ij})$).

For the D-Wave solver, we employ the Constrained Quadratic Model (CQM) formulation, which handles constraints natively without requiring penalty tuning.

\subsection{Computational Platforms}

Table~\ref{tab:platforms} summarizes the computational platforms used in this study.

\begin{table}[htbp]
\caption{Computational Platforms Summary}
\begin{center}
\begin{tabular}{|l|l|c|l|}
\hline
\textbf{Platform} & \textbf{Type} & \textbf{Qubits} & \textbf{Access} \\
\hline
Gurobi 13.0 & Classical MILP & -- & Local WLS \\
\hline
D-Wave Leap & Quantum Annealer & 5000+ & Cloud API \\
\hline
Rigetti Cepheus & Gate-based QPU & 108 & AWS Braket \\
\hline
Amazon SV1 & State Vector Sim & 34 & AWS Braket \\
\hline
Quanfluence & Ising Machine & 1000+ & REST API \\
\hline
\end{tabular}
\label{tab:platforms}
\end{center}
\end{table}

\subsubsection{Classical Solver}
We employ Gurobi Optimizer (version 13.0) with the Miller-Tucker-Zemlin (MTZ) formulation for subtour elimination. The solver runs on a workstation with an Intel Core i7 processor and 32GB RAM, using a Web License Server (WLS) configuration.

\subsubsection{D-Wave Quantum Annealer}
The D-Wave Leap hybrid Constrained Quadratic Model (CQM) solver is accessed via the Ocean SDK (version 6.x). This hybrid approach combines quantum annealing with classical pre- and post-processing, enabling the solution of problems larger than the native hardware connectivity would otherwise permit. The Advantage system features over 5,000 qubits with Pegasus topology connectivity.

\subsubsection{Gate-Based Quantum Computing}
QAOA circuits are executed on AWS Braket using both the SV1 state vector simulator and the Rigetti Cepheus-1-108Q superconducting quantum processor (108 qubits, located in us-west-1 region). Due to hardware gate count limitations (20,000 gates maximum on Rigetti), we employ a cluster-based decomposition strategy with maximum cluster sizes of 4 cities per QAOA sub-problem.

\subsubsection{Quantum-Inspired Ising Machine}
The Quanfluence Ising machine is accessed via REST API. This quantum-inspired optical computing platform solves QUBO formulations through coherent Ising machine dynamics, offering rapid solution times for moderately-sized optimization problems.

\subsection{QAOA Circuit Configuration}

The Quantum Approximate Optimization Algorithm (QAOA) alternates between a cost Hamiltonian $H_C$ encoding the QUBO objective and a mixer Hamiltonian $H_M = \sum_i X_i$ that enables state exploration. The parameterized circuit applies $p$ layers:
\begin{equation}
|\psi(\boldsymbol{\gamma}, \boldsymbol{\beta})\rangle = \prod_{l=1}^{p} e^{-i\beta_l H_M} e^{-i\gamma_l H_C} |+\rangle^{\otimes n}
\label{eq:qaoa}
\end{equation}

Table~\ref{tab:qaoa_params} details the QAOA parameters used for each hardware target.

\begin{table}[htbp]
\caption{QAOA Configuration Parameters}
\begin{center}
\begin{tabular}{|l|c|c|c|c|}
\hline
\textbf{Parameter} & \textbf{SV1} & \textbf{Rigetti} \\
\hline
QAOA layers ($p$) & 2 & 1 \\
\hline
Shots per circuit & 500 & 100 \\
\hline
Max cluster size & 5 cities & 4 cities \\
\hline
Max qubits/cluster & 25 & 16 \\
\hline
Est. gates/cluster & $\sim$10,000 & $\sim$4,300 \\
\hline
\end{tabular}
\label{tab:qaoa_params}
\end{center}
\end{table}

The gate count scales approximately as $O(n^4 \cdot p)$ for an $n$-city TSP QUBO due to the dense connectivity requirements. With Rigetti's 20,000 gate limit, we determined empirically that 4 cities ($n=4$, 16 qubits) with $p=1$ yields approximately 4,300 gates, providing adequate margin for reliable execution.

\subsection{Cluster-Based Decomposition}

To address scalability limitations inherent in current quantum hardware, we implement a hierarchical clustering approach uniformly across all quantum and quantum-inspired solvers:

\begin{enumerate}
\item \textbf{Coordinate Estimation}: For problems specified only by distance matrices, we employ Multidimensional Scaling (MDS) to recover approximate 2D city coordinates.
\item \textbf{K-Means Clustering}: Cities are partitioned into clusters using K-means, with cluster sizes bounded by hardware-specific limits (4 cities for QAOA, 10 for Quanfluence, 15 for D-Wave and Gurobi).
\item \textbf{Sub-Problem Solution}: Each cluster's internal TSP is solved independently using the respective quantum or classical solver.
\item \textbf{Tour Stitching}: Cluster tours are connected via a greedy nearest-neighbor heuristic based on cluster centroids.
\item \textbf{Local Search Refinement}: A 2-opt local search algorithm refines the stitched solution, typically yielding 20--40\% improvement in tour length.
\end{enumerate}

This decomposition enables scaling to 100+ city instances while maintaining the quantum advantage exploration within tractable sub-problems.

\subsection{Evaluation Metrics}

We evaluate solver performance using the following metrics:
\begin{itemize}
\item \textbf{Total Tour Distance}: The primary objective value representing the sum of edge costs in the solution tour.
\item \textbf{Solve Time}: Wall-clock time from problem submission to solution retrieval, including any queue delays for cloud-based quantum resources.
\item \textbf{Raw vs. Refined Distance}: For quantum solvers, we report both the initial solution quality and the post-2-opt refined distance.
\item \textbf{Cluster Count}: The number of sub-problems created by the decomposition strategy.
\item \textbf{Approximation Ratio}: Defined as $r = C_{\text{solver}} / C_{\text{optimal}}$, where $C_{\text{optimal}}$ is the best known solution (typically from Gurobi). A ratio of 1.0 indicates optimal solution quality.
\item \textbf{2-opt Improvement}: The percentage reduction in tour distance achieved by local search refinement, calculated as $(C_{\text{raw}} - C_{\text{refined}}) / C_{\text{raw}} \times 100\%$.
\end{itemize}

All experiments were conducted as single runs with fixed random seeds to ensure deterministic problem generation. For quantum hardware experiments, variability due to shot noise and hardware fluctuations is inherent; we report raw measurement outcomes without statistical averaging across multiple submissions.

\subsection{Reproducibility and Code Availability}

To facilitate reproducibility and enable independent verification of our results, all source code, benchmark scripts, and result data are publicly available. The repository includes:
\begin{itemize}
\item Solver implementations for all four platforms (Gurobi, D-Wave, QAOA, Quanfluence)
\item Cluster-based decomposition framework
\item Benchmark generation and execution scripts
\item Raw result data in JSON format
\end{itemize}

The codebase is available at: \texttt{https://github.com/SuperQ-Computing/universal-gurobi-controller}

Additionally, a web application is provided that allows readers to submit custom TD-TSP instances for solution by our implemented solvers, enabling hands-on verification of the reported results.

\section{Results}

This section presents the experimental results comparing classical, quantum, and quantum-inspired solvers across TSP instances of varying sizes. All solvers employ the cluster-based decomposition strategy described in Section IV to ensure fair comparison at scale.

\subsection{Solution Quality Comparison}

Table~\ref{tab:results_quality} summarizes the final tour distances achieved by each solver across all problem sizes. Bold values indicate the best (lowest) distance for each problem size.

\begin{table}[htbp]
\caption{Solution Quality: Final Tour Distance (units)}
\begin{center}
\begin{tabular}{|c|c|c|c|c|c|}
\hline
\textbf{Cities} & \textbf{Gurobi} & \textbf{D-Wave} & \textbf{Quanfl.} & \textbf{SV1} & \textbf{Rigetti} \\
\hline
5 & \textbf{227.35} & \textbf{227.35} & \textbf{227.35} & \textbf{227.35} & \textbf{227.35} \\
\hline
10 & \textbf{290.31} & \textbf{290.31} & \textbf{290.31} & \textbf{290.31} & \textbf{290.31} \\
\hline
25 & 418.58 & 415.96 & \textbf{407.30} & 449.06 & \textbf{407.30} \\
\hline
50 & 643.05 & 591.58 & \textbf{589.93} & 591.35 & 609.40 \\
\hline
100 & 824.33 & \textbf{781.56} & 934.98 & 782.62 & 790.13 \\
\hline
\end{tabular}
\label{tab:results_quality}
\end{center}
\end{table}

Key observations from the solution quality analysis:

\begin{itemize}
\item \textbf{Small instances (5--10 cities)}: All five solvers achieve identical optimal solutions, demonstrating that quantum and quantum-inspired approaches match classical optimality for tractable problem sizes.
\item \textbf{Medium instances (25--50 cities)}: The Quanfluence Ising machine produces the best solutions (407.30 and 589.93), outperforming even the clustered Gurobi approach. The Rigetti QPU matches Quanfluence on the 25-city instance.
\item \textbf{Large instance (100 cities)}: D-Wave achieves the best result (781.56), followed closely by QAOA on SV1 (782.62) and Rigetti (790.13). Quanfluence shows degraded performance at this scale (934.98).
\end{itemize}

\subsection{Computational Time Analysis}

Table~\ref{tab:results_time} presents the wall-clock solve times for each solver. These times include all overhead: QUBO construction, API communication, queue delays, and post-processing.

\begin{table}[htbp]
\caption{Computational Time (seconds)}
\begin{center}
\begin{tabular}{|c|c|c|c|c|c|}
\hline
\textbf{Cities} & \textbf{Gurobi} & \textbf{D-Wave} & \textbf{Quanfl.} & \textbf{SV1} & \textbf{Rigetti} \\
\hline
5 & 2.1 & 24.2 & \textbf{2.9} & 24.4 & $<$1 \\
\hline
10 & \textbf{0.3} & 20.8 & 2.9 & 13.9 & 19.5 \\
\hline
25 & \textbf{0.4} & 44.6 & 7.8 & 92.3 & 34.0 \\
\hline
50 & \textbf{0.6} & 85.4 & 8.4 & 89.2 & 92.9 \\
\hline
100 & \textbf{1.0} & 151.8 & 20.9 & 199.5 & 149.8 \\
\hline
\end{tabular}
\label{tab:results_time}
\end{center}
\end{table}

The classical Gurobi solver demonstrates superior speed across all problem sizes, completing the 100-city instance in approximately 1 second. Among quantum approaches, Quanfluence offers the fastest execution (20.9s for 100 cities), while QAOA on SV1 exhibits the longest runtimes due to the computational cost of state vector simulation. The Rigetti QPU shows comparable timing to D-Wave for larger instances.

\subsection{Real Quantum Hardware Validation}

A critical contribution of this work is the validation of our cluster-QAOA approach on real quantum hardware. Table~\ref{tab:qpu_comparison} compares results between the SV1 simulator and the Rigetti Cepheus-1-108Q QPU.

\begin{table}[htbp]
\caption{Simulator vs. Real QPU Comparison}
\begin{center}
\begin{tabular}{|c|c|c|c|}
\hline
\textbf{Cities} & \textbf{SV1 (Sim)} & \textbf{Rigetti (QPU)} & \textbf{Difference} \\
\hline
5 & 227.35 & 227.35 & 0.0\% \\
\hline
10 & 290.31 & 290.31 & 0.0\% \\
\hline
25 & 449.06 & 407.30 & --9.3\% \\
\hline
50 & 591.35 & 609.40 & +3.1\% \\
\hline
100 & 782.62 & 790.13 & +1.0\% \\
\hline
\end{tabular}
\label{tab:qpu_comparison}
\end{center}
\end{table}

The results demonstrate strong correlation between simulator and real hardware outcomes. For small instances, both achieve identical optimal solutions. For larger instances, the deviation remains within $\pm$10\%, with the QPU actually outperforming the simulator on the 25-city instance (--9.3\%). This validates that our cluster-based QAOA approach successfully transfers from simulation to real quantum hardware without significant degradation.

\subsection{Impact of 2-opt Post-Processing}

The 2-opt local search refinement proves essential for quantum solver performance. Table~\ref{tab:two_opt} shows the raw (pre-refinement) and final (post-refinement) distances for each quantum solver on the 100-city instance.

\begin{table}[htbp]
\caption{2-opt Improvement on 100-City Instance}
\begin{center}
\begin{tabular}{|l|c|c|c|}
\hline
\textbf{Solver} & \textbf{Raw Dist.} & \textbf{Final Dist.} & \textbf{Improvement} \\
\hline
D-Wave & 963.31 & 781.56 & 18.9\% \\
\hline
Quanfluence & 1555.78 & 934.98 & 39.9\% \\
\hline
QAOA (SV1) & 1224.44 & 782.62 & 36.1\% \\
\hline
QAOA (Rigetti) & 1140.60 & 790.13 & 30.7\% \\
\hline
\end{tabular}
\label{tab:two_opt}
\end{center}
\end{table}

The 2-opt refinement yields improvements ranging from 18.9\% (D-Wave) to 39.9\% (Quanfluence). Notably, D-Wave produces the highest-quality raw solutions, requiring the least post-processing. This suggests that D-Wave's hybrid CQM solver incorporates more sophisticated classical preprocessing, while pure QAOA and Ising machine approaches benefit substantially from local search enhancement.

\subsection{Approximation Ratio Analysis}

Using the best known solutions (from unclustered Gurobi where available) as baselines, we compute approximation ratios for each solver. Table~\ref{tab:approx_ratio} presents these ratios, where 1.00 indicates optimal and higher values indicate suboptimality.

\begin{table}[htbp]
\caption{Approximation Ratios (Lower is Better)}
\begin{center}
\begin{tabular}{|c|c|c|c|c|c|}
\hline
\textbf{Cities} & \textbf{Gurobi} & \textbf{D-Wave} & \textbf{Quanfl.} & \textbf{SV1} & \textbf{Rigetti} \\
\hline
5 & 1.00 & 1.00 & 1.00 & 1.00 & 1.00 \\
\hline
10 & 1.00 & 1.00 & 1.00 & 1.00 & 1.00 \\
\hline
25 & 1.03 & 1.02 & \textbf{1.00} & 1.10 & \textbf{1.00} \\
\hline
50 & 1.15 & 1.06 & \textbf{1.05} & 1.06 & 1.09 \\
\hline
100 & 1.05 & \textbf{1.00} & 1.20 & \textbf{1.00} & 1.01 \\
\hline
\end{tabular}
\label{tab:approx_ratio}
\end{center}
\end{table}

For the 100-city instance, D-Wave and QAOA (SV1) achieve optimal or near-optimal approximation ratios of 1.00, while clustered Gurobi shows a 5\% gap and Quanfluence exhibits a 20\% gap. This demonstrates that quantum approaches can match or exceed classical cluster-based heuristics for large-scale instances.

\subsection{Cluster Scaling Behavior}

The number of clusters required scales with problem size according to hardware-specific limits. Table~\ref{tab:clusters} shows the cluster counts for each solver.

\begin{table}[htbp]
\caption{Number of Clusters by Problem Size}
\begin{center}
\begin{tabular}{|c|c|c|c|c|}
\hline
\textbf{Cities} & \textbf{Gurobi} & \textbf{D-Wave} & \textbf{Quanfl.} & \textbf{QAOA} \\
\hline
5 & 1 & 1 & 1 & 1 \\
\hline
10 & 1 & 1 & 1 & 2 \\
\hline
25 & 2 & 2 & 3 & 6 \\
\hline
50 & 4 & 4 & 5 & 12 \\
\hline
100 & 7 & 7 & 10 & 25 \\
\hline
\end{tabular}
\label{tab:clusters}
\end{center}
\end{table}

QAOA requires the most clusters due to its stringent gate count constraints (max 4 cities per cluster), while Gurobi and D-Wave can handle larger clusters (up to 15 cities). Despite requiring more clusters, QAOA maintains competitive solution quality through effective 2-opt refinement.

\section{Discussion}

The experimental results reveal several important insights for the practical application of quantum computing to combinatorial optimization.

\subsection{Current State of Quantum Advantage}

Our benchmarks do not demonstrate quantum advantage in the traditional sense---Gurobi remains the fastest solver by a significant margin. However, quantum and quantum-inspired solvers achieve comparable or superior solution quality on medium-to-large instances when combined with classical post-processing. This suggests a hybrid paradigm where quantum resources contribute to solution exploration while classical algorithms refine results.

\subsection{Hardware-Specific Observations}

\textbf{D-Wave} excels at producing high-quality raw solutions, likely due to sophisticated classical preprocessing in the hybrid CQM solver. It achieves the best 100-city result despite moderate solve times.

\textbf{Quanfluence} offers the best speed-quality tradeoff for medium instances (25--50 cities) but struggles at 100 cities, possibly due to QUBO scaling limitations in the coherent Ising machine approach.

\textbf{QAOA on Rigetti} successfully validates the cluster-based approach on real quantum hardware, with results closely matching simulation. The 20,000 gate limit necessitates small clusters but does not prevent competitive performance.

\subsection{Importance of Hybrid Approaches}

The consistent 20--40\% improvement from 2-opt post-processing underscores that current quantum devices are best utilized within hybrid classical-quantum pipelines. Raw quantum solutions serve as high-quality starting points that classical local search can efficiently refine to near-optimality.

\subsection{Scalability Considerations}

The cluster-based decomposition enables scaling to arbitrarily large instances, but introduces a quality-scalability tradeoff. Smaller clusters (required for QAOA due to gate limits) produce more sub-optimal stitched tours, though 2-opt largely compensates. Future improvements in quantum hardware---higher gate counts, lower error rates---will enable larger clusters and potentially better raw solutions.

\section{Conclusion}

This paper presented a comprehensive benchmarking study comparing classical, quantum, and quantum-inspired approaches for solving the Time-Dependent Traveling Salesperson Problem. Our key contributions and findings are summarized below.

\subsection{Summary of Contributions}

\begin{enumerate}
\item \textbf{Multi-Platform Implementation}: We developed and evaluated TSP solvers across four distinct computational paradigms---classical MILP (Gurobi), quantum annealing (D-Wave), gate-based quantum computing (QAOA on AWS Braket), and quantum-inspired computing (Quanfluence Ising machine).

\item \textbf{Cluster-Based Scalability}: We introduced a unified cluster decomposition framework that enables all solvers to address problem sizes (up to 100 cities) beyond the native capabilities of current quantum hardware.

\item \textbf{Real QPU Validation}: We successfully executed QAOA on the Rigetti Cepheus-1-108Q quantum processor, demonstrating that cluster-based quantum optimization transfers effectively from simulation to real hardware with minimal performance degradation.

\item \textbf{Autonomous Code Generation}: We demonstrated the use of the Super platform for autonomous generation of quantum optimization code, with subsequent human refinement, establishing a workflow for accelerated quantum software development.
\end{enumerate}

\subsection{Key Findings}

Our experiments reveal that:
\begin{itemize}
\item All solvers achieve \textbf{identical optimal solutions} for small instances (5--10 cities), validating correctness across platforms.
\item \textbf{Quantum and quantum-inspired solvers match or exceed} clustered classical performance on medium-to-large instances when combined with 2-opt post-processing.
\item \textbf{D-Wave} produces the highest-quality raw solutions, achieving the best result on the 100-city benchmark.
\item \textbf{2-opt refinement is essential}, improving quantum solutions by 20--40\% and enabling competitive performance despite hardware limitations.
\item \textbf{Real QPU results correlate strongly} with simulation, with Rigetti achieving within 3\% of SV1 simulator performance.
\end{itemize}

\subsection{Future Work}

Several directions merit further investigation:
\begin{itemize}
\item \textbf{Larger problem scales}: Testing on instances with 500--1000+ cities to assess scaling behavior of the cluster approach.
\item \textbf{Alternative quantum algorithms}: Evaluating Grover-based approaches, quantum walks, and variational quantum linear solvers for TSP.
\item \textbf{Error mitigation}: Implementing zero-noise extrapolation and probabilistic error cancellation to improve real QPU solution quality.
\item \textbf{Time-dependent benchmarks}: Extending experiments to full TD-TSP with dynamic traffic patterns across multiple time windows.
\item \textbf{Hardware evolution tracking}: Repeating benchmarks as quantum hardware improves to document progress toward practical quantum advantage.
\end{itemize}

\subsection{Closing Remarks}

While current quantum hardware does not yet offer computational speedup over classical solvers for TSP, our results demonstrate that quantum approaches can achieve competitive solution quality within hybrid classical-quantum workflows. The cluster-based decomposition strategy provides a practical pathway for utilizing near-term quantum devices on real-world optimization problems. As quantum hardware continues to mature---with increasing qubit counts, improved gate fidelities, and higher gate count limits---we anticipate that the quality gap between raw quantum solutions and classical optima will narrow, potentially enabling genuine quantum advantage for combinatorial optimization in the coming years.

\section*{Acknowledgment}

The authors thank SuperQ Quantum Computing Inc. for providing access to the Super autonomous optimization platform and supporting this research. We acknowledge AWS for access to Amazon Braket quantum computing services, D-Wave Systems for access to the Leap quantum cloud platform, and Quanfluence for access to their Ising machine API. We also thank the University of Northern British Columbia for institutional support.

The computational experiments were conducted using Gurobi Optimizer under an academic license. The Rigetti Cepheus-1-108Q quantum processor was accessed through AWS Braket.

\section*{References}

Please number citations consecutively within brackets \cite{b1}. The 
sentence punctuation follows the bracket \cite{b2}. Refer simply to the reference 
number, as in \cite{b3}---do not use ``Ref. \cite{b3}'' or ``reference \cite{b3}'' except at 
the beginning of a sentence: ``Reference \cite{b3} was the first $\ldots$''

Number footnotes separately in superscripts. Place the actual footnote at 
the bottom of the column in which it was cited. Do not put footnotes in the 
abstract or reference list. Use letters for table footnotes.

Unless there are six authors or more give all authors' names; do not use 
``et al.''. Papers that have not been published, even if they have been 
submitted for publication, should be cited as ``unpublished'' \cite{b4}. Papers 
that have been accepted for publication should be cited as ``in press'' \cite{b5}. 
Capitalize only the first word in a paper title, except for proper nouns and 
element symbols.

For papers published in translation journals, please give the English 
citation first, followed by the original foreign-language citation \cite{b6}.

\begin{thebibliography}{00}
\bibitem{b1} G. Eason, B. Noble, and I. N. Sneddon, ``On certain integrals of Lipschitz-Hankel type involving products of Bessel functions,'' Phil. Trans. Roy. Soc. London, vol. A247, pp. 529--551, April 1955.
\bibitem{b2} J. Clerk Maxwell, A Treatise on Electricity and Magnetism, 3rd ed., vol. 2. Oxford: Clarendon, 1892, pp.68--73.
\bibitem{b3} I. S. Jacobs and C. P. Bean, ``Fine particles, thin films and exchange anisotropy,'' in Magnetism, vol. III, G. T. Rado and H. Suhl, Eds. New York: Academic, 1963, pp. 271--350.
\bibitem{b4} K. Elissa, ``Title of paper if known,'' unpublished.
\bibitem{b5} R. Nicole, ``Title of paper with only first word capitalized,'' J. Name Stand. Abbrev., in press.
\bibitem{b6} Y. Yorozu, M. Hirano, K. Oka, and Y. Tagawa, ``Electron spectroscopy studies on magneto-optical media and plastic substrate interface,'' IEEE Transl. J. Magn. Japan, vol. 2, pp. 740--741, August 1987 [Digests 9th Annual Conf. Magnetics Japan, p. 301, 1982].
\bibitem{b7} M. Young, The Technical Writer's Handbook. Mill Valley, CA: University Science, 1989.
\bibitem{b8} D. P. Kingma and M. Welling, ``Auto-encoding variational Bayes,'' 2013, arXiv:1312.6114. [Online]. Available: https://arxiv.org/abs/1312.6114
\bibitem{b9} S. Liu, ``Wi-Fi Energy Detection Testbed (12MTC),'' 2023, gitHub repository. [Online]. Available: https://github.com/liustone99/Wi-Fi-Energy-Detection-Testbed-12MTC
\bibitem{b10} ``Treatment episode data set: discharges (TEDS-D): concatenated, 2006 to 2009.'' U.S. Department of Health and Human Services, Substance Abuse and Mental Health Services Administration, Office of Applied Studies, August, 2013, DOI:10.3886/ICPSR30122.v2
\bibitem{b11} K. Eves and J. Valasek, ``Adaptive control for singularly perturbed systems examples,'' Code Ocean, Aug. 2023. [Online]. Available: https://codeocean.com/capsule/4989235/tree
\end{thebibliography}

\vspace{12pt}
\color{red}
IEEE conference templates contain guidance text for composing and formatting conference papers. Please ensure that all template text is removed from your conference paper prior to submission to the conference. Failure to remove the template text from your paper may result in your paper not being published.

\end{document}
